% ============================================================
% page.tex – Seitenstil, Fußnoten, Spacing-Helper
% ============================================================
% Diese Datei enthält feinabgestimmte Layout-Definitionen
% (Fußnoten, Microtype-Kerning, MS-Word-kompatibles Spacing).
% Eine Anpassung ist normalerweise nicht nötig.
% ============================================================


% Toleranz für Zeilenumbrüche lockern (vermeidet "Overfull hbox")
\sloppy


% ============================================================
% Kopf- und Fußzeile zurücksetzen
% ============================================================
\fancyhf{}
\rfoot{\thepage}
\renewcommand{\headrulewidth}{0pt}


% ============================================================
% Hilfsmittel
% ============================================================
% Zelle mit Zeilenumbrüchen in Tabellen:
%   \specialcell{Zeile 1\\ Zeile 2}
\newcommand{\specialcell}[2][c]{%
  \begin{tabular}[#1]{@{}t@{}}#2\end{tabular}}


% ============================================================
% Blockzitat-Umgebung
% ============================================================
% Eingerückt, einzeilig, kleinere Schrift – für wörtliche Zitate
% ab ca. drei Zeilen.
\renewenvironment{quotation}{
  \leftskip1cm
  \rightskip1cm
  \noindent
  \setstretch{1}
  \small
}


% ============================================================
% Fußnoten-Stil
% ============================================================
% Schmalere, dünne Trennlinie statt der Standard-Linie
\renewcommand\footnoterule{\kern-3pt \hrule width 3in height 0.7pt \hskip3pt \kern 2.6pt}

% 2 mm Einrückung nach der Fußnotennummer (SRH-Layout)
\let\oldfootnote\footnote
\renewcommand\footnote[1]{%
  \oldfootnote{\hspace{2mm}#1}}

% Fußnoten-Schriftgröße fest auf 10 pt
\renewcommand{\footnotesize}{\fontsize{10pt}{11pt}\selectfont}


% ============================================================
% Microtype-Feineinstellung
% ============================================================
% Zusätzliches Spacing um Gedankenstriche, Anführungszeichen und
% Klammern. Gilt nur für Helvetica/Arial fett (family={qhv}, series={b}).
\SetExtraKerning[unit=space]
  {
    encoding={*}, family={qhv}, series={b}, size={normalsize,large,Large}
  }
  {
    \textendash={400,400},        % Gedankenstrich
    \textquotedblleft={ ,120},    % öffnendes Anführungszeichen
    \textquotedblright={170, },   % schließendes Anführungszeichen
    "28={ ,250},                  % linke runde Klammer
    "29={300, }                   % rechte runde Klammer
  }


% ============================================================
% MS-Word-kompatible Zeilenabstände
% ============================================================
% Reproduziert MS Words "1,5-zeilig" und "doppelt" pro Schriftgröße.
% Verwendung im Dokument:
%   \MSonehalfspacing   1,5-zeilig
%   \MSdoublespacing    doppelt
%   \MSverbatimspacing  einzeilig (z.B. für Code/Verbatim)
\makeatletter
\newcommand{\MSonehalfspacing}{%
  \setstretch{1.44}%  Fallback
  \ifcase \@ptsize \relax % 10pt
    \setstretch {1.448}%
  \or % 11pt
    \setstretch {1.399}%
  \or % 12pt
    \setstretch {1.433}%
  \fi
}
\newcommand{\MSdoublespacing}{%
  \setstretch {1.92}%  Fallback
  \ifcase \@ptsize \relax % 10pt
    \setstretch {1.936}%
  \or % 11pt
    \setstretch {1.866}%
  \or % 12pt
    \setstretch {1.902}%
  \fi
}
\newcommand{\MSverbatimspacing}{%
  \setstretch{1.44}%  Fallback
  \ifcase \@ptsize \relax % 10pt
    \setstretch {1}%
  \or % 11pt
    \setstretch {1}%
  \or % 12pt
    \setstretch {1}%
  \fi
}
\makeatother
